\section{Privacy}
\label{sec:privacy}

\Zoo{} \DEX{} ships two complementary privacy paths via the \Lux{}
chain topology: optional \FHE{}-encrypted orderflow via F-Chain
\cite{lp-013}, and \ZKP{}-based selective disclosure via Z-Chain
\cite{lp-063}.

\subsection{F-Chain FHE for orderflow}

F-Chain (FVM) is a fully-homomorphic-encryption execution environment
\cite{lp-013}. \Zoo{} \DEX{} contracts can dispatch orderflow to
F-Chain so that orders, balances, and fills are ciphertexts. A
confidential trade looks like:

\begin{enumerate}[leftmargin=1.4em]
  \item User's wallet encrypts the order under the F-Chain public
        key. The cleartext (price, quantity, side) never leaves the
        wallet.
  \item The F-Chain contract receives the ciphertext order and runs
        the matching engine over the ciphertexts: comparison,
        crossing detection, and fill computation are all
        ciphertext operations.
  \item The fill receipt is returned to the user as a ciphertext;
        the user decrypts locally to see the fill price and quantity.
  \item The trade is anchored to D-Chain with a public commitment
        (the encrypted balance update), but no observer can read
        the cleartext order or fill.
\end{enumerate}

Use F-Chain confidential orderflow when:

\begin{itemize}[leftmargin=1.2em]
  \item Front-running risk is material (illiquid securities,
        institutional block flow, accumulation strategies).
  \item Privacy of the position itself is part of the user's
        compliance posture (treasury holdings disclosed only at
        regulator-defined intervals).
  \item Per-\LLM{} token holdings should not telegraph governance
        positioning.
\end{itemize}

\paragraph{Confidential balances (LP-067).}
The position itself can live as a Confidential ERC-20 balance.
Holders' balances are ciphertexts; transfers update ciphertexts
without revealing amounts. The \Liquidity{} precompile can verify
holder eligibility on the cleartext side via TEE-attested decryption
where regulator obligations require it (e.g., transfer-agent updates
need cleartext beneficial-owner identification, but the publicly
visible chain log only carries ciphertexts).

\subsection{Z-Chain ZKP for selective disclosure}

Z-Chain (ZVM) is a zero-knowledge proof execution and registry
environment \cite{lp-063}. \Zoo{} \DEX{} dispatches verification
jobs to Z-Chain via precompile.

\paragraph{Accreditation proofs.}
A user proves ``I am an accredited investor'' without revealing
income, net worth, or jurisdiction. The proof is generated against
a circuit that takes the user's \Liquidity{}-issued accreditation
attestation as a private witness and outputs a public ``yes''.

\paragraph{Geographic eligibility.}
A user proves ``I am not in a sanctioned jurisdiction'' without
revealing the actual jurisdiction. The proof is generated against
a circuit that takes a geographic-attestation private witness and
outputs the negative-list result.

\paragraph{Source-of-funds.}
A user proves ``these funds came from a regulated source''
without revealing the source identity or transaction history. Used
for AML compliance on large positions.

\paragraph{Rolled-up KYC.}
For high-volume venues, Z-Chain rolls up batches of KYC eligibility
checks into a single Groth16 verification. The matching engine
consults the rolled-up proof rather than per-order proof
verification.

\subsection{When to use which}

\begin{longtable}{p{0.30\textwidth}p{0.30\textwidth}p{0.34\textwidth}}
\toprule
\textbf{Property to protect} & \textbf{F-Chain (\FHE{})} & \textbf{Z-Chain (\ZKP{})} \\
\midrule
Hide order price/qty & Yes & No \\
Hide balance & Yes (LP-067) & No (only the witness) \\
Prove accreditation & Indirect (via TEE) & Direct (predicate) \\
Hide identity & Via key management & Yes (predicate over identity) \\
Per-op cost & High & Low to moderate \\
Best for & Confidential orderflow, balances, block flow & Eligibility predicates, regulatory checks \\
\bottomrule
\end{longtable}

The wallet exposes a privacy-mode toggle (\S\ref{sec:retail-ux})
that picks the right chain for the user's intent. ``Hide this
trade'' routes to F-Chain confidential orderflow. ``Prove
eligibility'' routes to Z-Chain selective disclosure.

\subsection{Privacy and compliance compose}

A common confusion is whether privacy and compliance are at odds.
On \Zoo{} \DEX{} they compose:

\begin{itemize}[leftmargin=1.2em]
  \item The \Liquidity{} precompile sees what it needs to see for
        regulatory obligations (KYC status, accreditation, transfer
        eligibility) but does not see public order data.
  \item Public observers see what they need to see for fair price
        discovery (BBO, last trade) but do not see the user's
        identity or holdings.
  \item Regulators see the full audit trail through the \Liquidity{}
        Protocol's audit interface, with appropriate authorisation,
        but do not have a public window into the order book.
\end{itemize}

The compliance perimeter is at the \Liquidity{} precompile; the
privacy perimeter is at the F-Chain / Z-Chain boundary. They are
orthogonal axes; both can be enabled simultaneously.

\subsection{GPU-native privacy}

Both privacy paths are \GPU{}-native end-to-end. F-Chain TFHE
bootstrapping runs on the validator \GPU{} per LP-013; Z-Chain
Groth16 verification runs on the \GPU{} per LP-063. The privacy
primitives in \Zoo{} \DEX{} are not bolted on at application
layer; they execute on the same substrate as the rest of the chain.

\subsection{Key management for privacy mode}

\begin{itemize}[leftmargin=1.2em]
  \item \textbf{F-Chain encryption keys.} Generated client-side;
        the user holds the decryption key in the wallet. Loss of
        the wallet means loss of the private balance unless the
        user has set up social recovery via the \Lux{} Wallet
        recovery flow.
  \item \textbf{Z-Chain proving keys.} Public-circuit-specific.
        No user-held secret; the proof generator uses the user's
        \Liquidity{}-issued attestation as the witness.
  \item \textbf{Wallet privacy mode toggle.} Per-session, with
        explicit confirmation when toggling to confidential mode
        for a regulated security (because the wallet warns about
        the lookup-cost trade-off).
\end{itemize}

\subsection{Audit access}

For regulator-authorised audit, the \Liquidity{} Protocol audit
interface (\cite{liquidity-protocol-formal}) provides a controlled
path for cleartext access to the records that matter. The chain log
itself does not retain decryption keys; audit access is per-record,
authorisation-gated, and logged on chain.
